\documentclass[10pt,a4paper]{amsart}
\usepackage{amsmath,amssymb,amsthm,graphicx}
% Retained minimal preamble for the Stacks Project chapter more-algebra.tex (More on Algebra).
% The chapter file pinned-original-source.tex carries only \input{preamble} on its first line, so no
% preamble text lives inside the chapter itself.  The Stacks Project uses one single shared preamble for
% every chapter; the snapshot stages it identically as sources/stacks-chow/preamble.tex and
% sources/stacks-etale-cohomology/preamble.tex.  The block below is copied mechanically, line for line,
% from sources/stacks-chow/preamble.tex lines 14-18 (the Xy-pic setup) and 154-217 (the native theorem
% declarations together with the mathematical macros).  Nothing was added, removed or redefined; the
% wrapper class and the amsmath/amssymb/amsthm/graphicx base packages are supplied by the assembler.
% Chapter: more-algebra.tex -- More on Algebra.
% For commutative diagrams we use Xy-pic
\usepackage[all]{xy}

% We use 2cell for 2-commutative diagrams.
\xyoption{2cell}

% Theorem environments.
%
\theoremstyle{plain}
\newtheorem{theorem}[subsection]{Theorem}
\newtheorem{proposition}[subsection]{Proposition}
\newtheorem{lemma}[subsection]{Lemma}

\theoremstyle{definition}
\newtheorem{definition}[subsection]{Definition}
\newtheorem{example}[subsection]{Example}
\newtheorem{exercise}[subsection]{Exercise}
\newtheorem{situation}[subsection]{Situation}

\theoremstyle{remark}
\newtheorem{remark}[subsection]{Remark}
\newtheorem{remarks}[subsection]{Remarks}

\numberwithin{equation}{subsection}

% Macros
%
\def\lim{\mathop{\mathrm{lim}}\nolimits}
\def\colim{\mathop{\mathrm{colim}}\nolimits}
\def\Spec{\mathop{\mathrm{Spec}}}
\def\Hom{\mathop{\mathrm{Hom}}\nolimits}
\def\Ext{\mathop{\mathrm{Ext}}\nolimits}
\def\SheafHom{\mathop{\mathcal{H}\!\mathit{om}}\nolimits}
\def\SheafExt{\mathop{\mathcal{E}\!\mathit{xt}}\nolimits}
\def\Sch{\mathit{Sch}}
\def\Mor{\mathop{\mathrm{Mor}}\nolimits}
\def\Ob{\mathop{\mathrm{Ob}}\nolimits}
\def\Sh{\mathop{\mathit{Sh}}\nolimits}
\def\NL{\mathop{N\!L}\nolimits}
\def\CH{\mathop{\mathrm{CH}}\nolimits}
\def\proetale{{pro\text{-}\acute{e}tale}}
\def\etale{{\acute{e}tale}}
\def\QCoh{\mathit{QCoh}}
\def\Ker{\mathop{\mathrm{Ker}}}
\def\Im{\mathop{\mathrm{Im}}}
\def\Coker{\mathop{\mathrm{Coker}}}
\def\Coim{\mathop{\mathrm{Coim}}}

% Boxtimes
%
\DeclareMathSymbol{\boxtimes}{\mathbin}{AMSa}{"02}

%
% Macros for moduli stacks/spaces
%
\def\QCohstack{\mathcal{QC}\!\mathit{oh}}
\def\Cohstack{\mathcal{C}\!\mathit{oh}}
\def\Spacesstack{\mathcal{S}\!\mathit{paces}}
\def\Quotfunctor{\mathrm{Quot}}
\def\Hilbfunctor{\mathrm{Hilb}}
\def\Curvesstack{\mathcal{C}\!\mathit{urves}}
\def\Polarizedstack{\mathcal{P}\!\mathit{olarized}}
\def\Complexesstack{\mathcal{C}\!\mathit{omplexes}}
% \Pic is the operator that assigns to X its picard group, usage \Pic(X)
% \Picardstack_{X/B} denotes the Picard stack of X over B
% \Picardfunctor_{X/B} denotes the Picard functor of X over B
\def\Pic{\mathop{\mathrm{Pic}}\nolimits}
\def\Picardstack{\mathcal{P}\!\mathit{ic}}
\def\Picardfunctor{\mathrm{Pic}}
\def\Deformationcategory{\mathcal{D}\!\mathit{ef}}
\usepackage{cleveref}
\crefname{originallemma}{original lemma}{original lemmas}
\Crefname{originallemma}{Original lemma}{Original lemmas}
\makeatletter
\newlabel{lemma-vanishing-beta}{{0F7Y}{source}{}{sourcepointer-1}{}}
\newlabel{lemma-vanishing-beta@cref}{{[originallemma][0][]0F7Y}{[1][0][]source}{}{}{}}
\makeatother
\crefname{originallemma}{original lemma}{original lemmas}
\Crefname{originallemma}{Original lemma}{Original lemmas}
\makeatletter
\newlabel{lemma-eta-vanishing-beta}{{0F7Z}{source}{}{sourcepointer-2}{}}
\newlabel{lemma-eta-vanishing-beta@cref}{{[originallemma][0][]0F7Z}{[1][0][]source}{}{}{}}
\makeatother
\crefname{originallemma}{original lemma}{original lemmas}
\Crefname{originallemma}{Original lemma}{Original lemmas}
\makeatletter
\newlabel{lemma-eta-locally-free}{{0F7W}{source}{}{sourcepointer-3}{}}
\newlabel{lemma-eta-locally-free@cref}{{[originallemma][0][]0F7W}{[1][0][]source}{}{}{}}
\makeatother
\title{The ideal of the Bockstein operators is finitely generated and its quotient makes the shifted complex have locally free cohomology}
\author{The Stacks Project Authors}
\date{stacks-project:more-algebra.tex:local-snapshot-20261012; GNU FDL 1.2 or later}
\begin{document}
\maketitle
\setcounter{subsection}{1}
\noindent\textit{Editorial note.} All delivered mathematical text is the author's own, copied from the pinned Stacks Project chapter file more-algebra.tex; no statement or proof step was rewritten, reordered or removed. The counted statement is the labeled lemma environment at source lines 28197-28228 (canonical tag 0F81) and the counted proof is the proof environment at source lines 28230-28306, which immediately follows it and is its own proof.\par\medskip
\noindent\textit{Editorial note.} The statement uses the letter f both for the fixed nonzerodivisor and inside the displayed set-builder notation for primes containing it, and it names the sum of the ideals J\_i as J of M and f; both are printed as the author wrote them.\par\medskip
\noindent\textit{Editorial note.} The proof leaves two short empty stretches of formatting at source lines 28261-28263 between the proof of (3) and the proof of (5); they are part of the counted proof span and are delivered as they stand.\par\medskip
\noindent\textit{Editorial note.} The proof refers to three earlier lemmas of the same chapter, namely the two vanishing beta statements and the local freeness of the shifted complex; all three are carried as navigation stubs with their roles recorded in the reasons.\par\medskip
\noindent\textit{Editorial note.} The chapter file pinned-original-source.tex carries only an input of the shared preamble on its first line, so no preamble text lives inside the chapter. The retained file sources/stacks-more-algebra/retained-preamble.tex is a mechanical, line-for-line copy of the project's single shared preamble as staged in the snapshot (Xy-pic setup plus the native theorem declarations and mathematical macros). Nothing was added, removed or redefined, and the excerpt serves only as preamble, never as a span coordinate.\par\medskip
\noindent\textit{Editorial note.} The source URL is the source-level Stacks Project URL, as required when more than one card is produced from one source; the per-claim tag is carried by claim\_key and by this note. The counted statement and proof cite no bibliography entry, so this card needs no bibliography.\par\medskip
\noindent\textit{Allowed standard prerequisites (editorial).}\begin{itemize}
\item The Bockstein operators beta attached to a bounded complex of finite free modules and a nonzerodivisor vanish exactly where the reduction modulo f\textasciicircum{}2 surjects onto the reduction modulo f, and their vanishing implies the condition defining the locus E. used at proof lines 28231-28233 and again at proof lines 28250-28252. Identifies the two descriptions of E and supplies part (4).
\item The shifted complex has finite locally free terms of the given ranks when the Fitting ideals are principal. used at proof lines 28241-28242. Makes the terms of the shifted complex finite locally free, which is what allows the local splitting in the proof of part (3).
\item An ideal defined by a universal property of a map between free modules contains exactly the elements by which the map can be made to split as prescribed, so the local vanishing of that ideal is read off from the existence of a compatible direct sum decomposition. used at proof lines 28256-28259. Produces the vanishing of the localized ideals J\_i from the displayed splitting, which is the whole content of part (3).
\end{itemize}\medskip
\ifdefined\hypertarget\hypertarget{sourcepointer-1}{}\fi
\noindent\textit{Source reference.} Original lemma 0F7Y: sources/stacks-more-algebra/pinned-original-source.tex, source lines 27783--27783. The reference is invoked at proof lines 28231-28233 to equate the two descriptions of E and to supply part (4); the proofs of the remaining parts are inside the counted span.\par\medskip
\ifdefined\hypertarget\hypertarget{sourcepointer-2}{}\fi
\noindent\textit{Source reference.} Original lemma 0F7Z: sources/stacks-more-algebra/pinned-original-source.tex, source lines 27807--27807. The reference is used at proof lines 28250-28252 to obtain the compatible direct sum decomposition of the localized reduction, which is the input to the universal property argument in part (3); that argument is inside the counted span.\par\medskip
\ifdefined\hypertarget\hypertarget{sourcepointer-3}{}\fi
\noindent\textit{Source reference.} Original lemma 0F7W: sources/stacks-more-algebra/pinned-original-source.tex, source lines 28050--28050. The reference gives the finite local freeness of the terms of the shifted complex, used at proof lines 28241-28242 as the reason the localized reductions split; everything done with that splitting is inside the counted span.\par\medskip
\begin{lemma}
\label{lemma-eta-vanishing-beta-plus}
Let $A$ be a ring and let $f \in A$ be a nonzerodivisor. Let
$M^\bullet$ be a bounded complex of finite free $A$-modules.
Assume $I_i(M^\bullet, f)$ is a principal ideal
for all $i \in \mathbf{Z}$. Consider the ideal
$J(M^\bullet, f) = \sum_i J_i(M^\bullet, f)$ of $A/fA$.
Consider the set of prime ideals
\begin{align*}
E
& =
\{f \in \mathfrak p \subset A \mid
\Ker(d^i \bmod f^2)_\mathfrak p \text{ surjects onto }
\Ker(d^i \bmod f)_\mathfrak p \text{ for all }i \in \mathbf{Z}\} \\
& =
\{f \in \mathfrak p \subset A \mid
\text{the localizations }\beta_\mathfrak p
\text{ of the Bockstein operators are zero}\}
\end{align*}
Then we have
\begin{enumerate}
\item $J(M^\bullet, f)$ is finitely generated,
\item $A/fA \to C = (A/fA)/J(M^\bullet, f)$
is surjective of finite presentation,
\item $J(M^\bullet, f)_\mathfrak p = 0$ for $\mathfrak p \in E$,
\item if $f \in \mathfrak p$ and
$H^i(M^\bullet)_\mathfrak p$ is free for all $i \in \mathbf{Z}$,
then $\mathfrak p \in E$, and
\item the cohomology modules of
$\eta_f M^\bullet \otimes_A C$ are finite locally free $C$-modules.
\end{enumerate}
\end{lemma}

\begin{proof}
The equality in the definition of $E$ follows from
Lemma \ref{lemma-vanishing-beta} and in addition the
final statement of that lemma implies part (4).

\medskip\noindent
Part (1) is true because the ideals $J_i(M^\bullet, f)$ are finitely
generated and because $M^\bullet$ is bounded and hence $J_i(M^\bullet, f)$
is zero for almost all $i$. Part (2) is just a reformulation of part (1).

\medskip\noindent
Proof of (3). By Lemma \ref{lemma-eta-locally-free} we find that
$(\eta_fM)^i$ is finite locally free of rank $r_i$ for all $i$.
Consider the map
$$
(1, d^i) :
(\eta_fM)^i / f(\eta_fM)^i
\longrightarrow
f^iM^i/f^{i + 1}M^i \oplus f^{i + 1}M^{i + 1}/f^{i + 2}M^{i + 1}
$$
Pick $\mathfrak p \in E$. By Lemma \ref{lemma-eta-vanishing-beta}
and the local freeness of the modules $(\eta_fM)^i$ we may write
$$
\left((\eta_fM)^i / f(\eta_fM)^i\right)_\mathfrak p =
(A/fA)_\mathfrak p^{\oplus m_i} \oplus (A/fA)_\mathfrak p^{\oplus n_i}
$$
compatible with the arrow $(1, d^i)$ above. By the universal property of
the ideal $J_i(M^\bullet, f)$ we conclude that
$J_i(M^\bullet, f)_\mathfrak p = 0$.
Hence $I_\mathfrak p = fA_\mathfrak p$ for $\mathfrak p \in E$.

\medskip\noindent


\medskip\noindent
Proof of (5). Observe that the differential on $\eta_fM^\bullet$
fits into a commutative diagram
$$
\xymatrix{
(\eta_fM)^i \ar[d] \ar[r] &
f^iM^i \oplus f^{i + 1}M^{i + 1}
\ar[d]^{\left(
\begin{matrix}
0 & 1 \\
0 & 0
\end{matrix}
\right)} \\
(\eta_fM)^{i + 1} \ar[r] &
f^{i + 1}M^i \oplus f^{i + 2}M^{i + 2}
}
$$
By construction, after tensoring with $C$, the modules on the
left are direct sums of direct summands of the summands on the right.
Picture
$$
\xymatrix{
(\eta_fM)^i \otimes_A C \ar[d] \ar@{=}[r] &
K^i \oplus L^i \ar[r] \ar[d] &
f^iM^i \otimes_A C \oplus f^{i + 1}M^{i + 1} \otimes_A C
\ar[d]^{\left(
\begin{matrix}
0 & 1 \\
0 & 0
\end{matrix}
\right)} \\
(\eta_fM)^{i + 1} \otimes_A C \ar@{=}[r] &
K^{i + 1} \oplus L^{i + 1} \ar[r] &
f^{i + 1}M^i \otimes_A C \oplus f^{i + 2}M^{i + 2} \otimes_A C
}
$$
where the horizontal arrows are compatible with direct sum decompositions
as well as inclusions of direct summands.
It follows that the differential identifies $L^i$ with a direct summand
of $K^{i + 1}$ and we conclude that the cohomology of
$\eta_fM^\bullet \otimes_A C$ in degree $i$ is the
module $K^{i + 1}/L^i$ which is finite projective as desired.
\end{proof}

\end{document}
